\chapter{Арга зүй}

\section{Системийн шаардлага}
Энд системийн функциональ болон функциональ бус шаардлагуудыг тодорхойлсон.

\subsection{Функциональ шаардлагууд}

\begin{itemize}
	\item ФШ-01 --- R2 робот нь жадны үзүүр авч, R1 роботтой хамт зэвсгээ угсрах шаардлагатай.
	\item ФШ-02 --- R2 робот нь мэйхуа ойг даван туулах бөгөөд, дор хаяж нэг судар авч гарах шаардлагатай.
	\item ФШ-03 --- R2 робот нь икс окс тавиурт авч ирсэн судараа байрлуулах шаардлагатай.
\end{itemize}

\subsection{Функциональ бус шаардлагууд}

\subsubsection{Программ хангамжийн шаардлагууд}
\begin{itemize}
	\item ФБШ-ПХ-01 --- R2 робот нь секундэд 10 удаа камераар авсан датаг шүүж, камерт үзэгдэх обьектуудын байршлыг тогтоох шаардлагатай.
	\item ФБШ-ПХ-02 --- R2 робот нь өөрийн байршлыг таниж, гүйцэтгэж буй явцаа ойлгосон байх шаардлагатай.
	\item ФБШ-ПХ-03 --- R2 робот нь бүрэн автоматаар ажиллах шаардлагатай.
\end{itemize}

\subsubsection{Гүйцэтгэлийн шаардлагууд}
\begin{itemize}
	\item ФБШ-Г-01 --- R2 робот нь зэвсэг угсрахдаа R1 роботтой шүргэлцэхгүй байх шаардлагатай.
	\item ФБШ-Г-02 --- R2 робот нь мэйхуа ойг давж гарах чадвартай байх шаардлагатай.
	\item ФБШ-Г-03 --- R2 робот нь судар авах, тавих болон аваад явахдаа унагахгүй байх шаардлагатай.
\end{itemize}

\newpage
\section{Системийн архитектур}

\subsection{Ерөнхий системийн архитектур}

Энэ хэсэгт ерөнхий системийн архитектур болох дүрс таних, шийдвэр гаргах гэсэн үндсэн 2 дэд систем болон тэдгээрийн R2 роботтой холбогдох схемийг илэрхийлсэн болно.

\begin{figure}[h]
    \centering
    \includegraphics[width=\textwidth]{robot-structure.drawio.png}
    \caption{Ерөнхий системийн архитектур}
    \label{fig:robot-structure}
\end{figure}

\subsection{Датасет боловсруулах системийн архитектур}

Энэ хэсэгт обьект таних алгоритм сургахын тулд хэрэглэгдэх датасет боловсруулах системийн архитектурийг илэрхийлсэн болно.

\begin{figure}[h]
    \centering
    \includegraphics[width=0.7\textwidth]{dataset-generation-architecture.drawio.png}
    \caption{Датасет боловсруулах системийн архитектур}
    \label{fig:dataset-generation-architecture}
\end{figure}

\newpage
\subsection{Шийдвэр гаргах алгоритмын төлөвийн диаграм}

Шийдвэр гаргах алгоритмыг төлөвийн автоматын дагуу загварчилсан болно.
Төлөвийн диаграм нь ерөнхий 4 хэсэгт хуваагдсан болно:

\begin{enumerate}
	\item \textbf{Тохиргоо} --- Баг болон сударын байршлуудыг тохируулах хэсэг. Энд тоглогч нар гараар оролт оруулж тохируулахаар байна.
	\item \textbf{Зэвсэг угсрах} --- Жадны үзүүр олж авч, R1 робот очиж угсартал зогсож хүлээж байна. Угсарсны дараа R1 робот салгаж авах бөгөөд, үүнийг мэдрэгч ашиглан мэдэрснээр дараагийн төлөвт шилжинэ.
	\item \textbf{Судар авах} --- Мэйхуа ойруу орж, 1 ширхэг судар авч гарна. Орох болон гарахдаа R2-ийн орох болон гарах хэсгээс хэтрэхгүй байхыг анхаарах хэрэгтэй.
	\item \textbf{Судар байршуулах} --- Тулааны талбайруу орж, хоосон нүд олоод судар байршуулна. Үүний дараа дахин давтагдах бөгөөд хэрэв жадны үзүүр олсон тохиолдолд зэвсэг угсрах хэсэгт шилжих бол олдоогүй тохиолдолд судар авах хэсэгт шилжинэ.
\end{enumerate}

\begin{figure}[h]
    \centering
    \includegraphics[width=0.7\textwidth]{state-diagram.drawio.png}
    \caption{Шийдвэр гаргах алгоритмын төлөвийн диаграм}
    \label{fig:decision-state-diagram}
\end{figure}

\newpage
\section{Обьект таних алгоритм сургалт}

R2 робот Мейхуан ойгоор явахдаа дор хаяж нэг судар авах хэрэгтэй бөгөөд, энэ судар нь икс окс тоглох гол зорилгыг биелүүлж өгдөг тул хамгийн чухал хэсэг болж өгсөн.

\subsection{Симуляцийн орчин бүрдүүлэх}

\subsubsection{Тоглох талбайн модель үүсгэх}

Датасет үүсгэхэд зориулагдсан симуляцийн орчинг бүрдүүлэх зорилготойгоор тоглох талбайн моделийг загварчилсан.
Энэхүү моделийн зураг төслийг \Robocon{} тэмцээний дүрэмд хавсаргасан зурагнаас сэдэвлэн үүсгэсэн \cite{robocon2026rulebook}.

\begin{figure}[h]
    \centering
    \includegraphics[width=0.7\textwidth]{robocon-field-blueprint.png}
    \caption{\Robocon{} тэмцээний дүрэмд хавсаргасан зураг төсөл \cite{robocon2026rulebook}}
    \label{fig:field-blueprint}
\end{figure}

Үүнийг загварчлахад FreeCAD болон Blender ашигласан.
FreeCAD нь геометрийн параметрууд ашиглан загварчилдаг бөгөөд, Робокон тэмцээний дүрэмд хавсаргасан тоглох талбайн зураг төслийг хавсаргасан байсан бөгөөд, үүний дагуу моделоо үүсгэсэн болно.

Моделоо үүсгэсний дараа Blender ашиглаж, өнгө болон материал өгсөн бөгөөд, энэ нь цаашлан Unity-руу моделоо оруулахад хэрэглэсэн болно.

\begin{figure}[h]
    \centering
    \begin{minipage}{0.45\textwidth}
        \includegraphics[width=\linewidth]{field-model-freecad.png}
        \caption{\Robocon{} тэмцээний тоглох талбайн модель}
        \label{fig:field-model-freecad}
    \end{minipage}
    \hfill
    \begin{minipage}{0.45\textwidth}
        \includegraphics[width=\linewidth]{field-model-blender.png}
        \caption{Материал олгосны дараах \Robocon{} тэмцээний тоглох талбайн модель}
        \label{fig:field-model-blender}
    \end{minipage}
\end{figure}

\newpage
\subsection{Обьект таниулах датасет үүсгэх}

Датасет үүсгэхэд дор хаяж 100 зураг хэрэгтэй бөгөөд, их байх тусам үр дүнтэй байдаг тул симуляцийн орчинд үүсгэхээр Unity видео тоглоомын хөдөлгүүрийг ашигласан болно.

\subsubsection{Датасет үүсгэх код}

Датасет үүсгэхэд төрөл бүрийн өнцгөөс сударын зургуудыг авах хэрэгтэй болсон. Үүнд:
\begin{enumerate}
    \item Тоглогчийн хөдөлгөөний код
    \item Камер байршуулах код
    \item Датасет үүсгэх код
\end{enumerate}
хэмээн 3 ангилж бичсэн.

\subsubsection{Тоглогчийн хөдөлгөөний код}

Энэ нь энгийн бөгөөд тухайн тоглогч өөрийн байршлаа сольж, зураг авах тохиромжтой байрлалд очиход зориулагдсан билээ.

\subsubsection{Камер байршуулах код}

Датасет үүсгэхэд сударуудыг төрөл бүрийн өнцгөөс авсан зураг шаардлагатай бөгөөд, үүнд зориулсан камерийн дурын байршилд байрлуулах код бичсэн юм.
Энэ нь тоглогчийн байршлаас өгөгдсөн хүрээнд дурын зай, өнцөгт байрлах бөгөөд, зураг болгон дурын өнцөгт байрлуулсан байх нь үр дүнг нэмэгдүүлдэг.

\subsubsection{Датасет үүсгэх код}

Датасет үүсгэх код нь камер байршуулах код камер байршуулах бүрт ажиллах бөгөөд, энэ нь камерийн дүрсийг файлд хадгалж, мөн дүрсэд агуулагдах, сударуудын байрлалыг тооцоолж, камерт ямархуу зай эзэлж буйг тооцоолон, хариултуудыг үүсгэдэг.

Симуляцийн орчинд судар болгоны байршил нь тоглоомын ертөнцийн зүгээс өгөгдсөн байдаг тул үүнийг камерийн зүгээс харагдах байршил болгон хувиргах тооцоолол хийдэг.
